% Generated from locale/tr/content/set-theory/replacement/absinf.tex; do not hand-edit.


\olfileid[tr]{sth}{replacement}{absinf}
\olsection{Yerine Koyma ve ``Mutlak Sonsuzluk''}

Şimdi \limofsize{} ilkesini geride bırakıp Yerine Koyma'yı
gerekçelendirmeye yönelik, kümelerin aşamalar hâlinde oluşturulması fikrini
ciddiye alan farklı bir (içsel) girişimler ailesini inceleyeceğiz.

Yinelemeli süreci ilk kez ana hatlarıyla ortaya koyduğumuzda her aşamada ne
olduğunu açıklayan bazı ilkeler sunduk. Bunlar \stageshier, \stagesord{} ve
\stagesacc{} idi. Daha sonra aşamaların sayısı hakkında bir şeyler söyleyen
bazı ilkeler ekledik: \stagessucc{} bize küme oluşturma sürecinin hiçbir zaman
sona ermediğini, \stagesinf{} ise sürecin sonsuzuncu bir aşamadan geçtiğini
söyledi.

Bu son iki ilkenin şöyle daha geniş bir ilkeden çıktığını öne sürmek
akla yatkındır:
\begin{enumerate}
	\item[] \stagesinex. Mutlak olarak sonsuz sayıda aşama vardır;
	hiyerarşi olabileceği kadar yüksektir.
\end{enumerate}
Bu elbette gayriresmî bir ilkedir. Fakat birikimli-yinelemeli küme anlayışı
tarafından doğrudan \emph{gerektirilmese} bile onunla kesinlikle
\emph{bağdaşır} görünür. En azından \limofsize{} ilkesinin aksine, kümelerin
aşama aşama oluşturulduğu fikrini korur.

Şimdiki umut, \stagesinex{} ilkesinden yararlanarak Yerine Koyma'yı
gerekçelendirmektir. Bunun nasıl yapılabileceğine bakalım.

\olref[ordinals][idea]{sec} içinde, (özü itibarıyla) $\omega+\omega$'ya
izomorf olması gereken bir iyi sıralama kurmanın kolay olduğunu gördük. Başka
deyişle sıra tipi (özü itibarıyla) $\omega+\omega$ olan, aşama aşama ilerleyen
bir yinelemeli süreci kolayca tasarlayabiliriz. Dolayısıyla \stagesinex{}
ilkesini kabul etmişsek hiyerarşide en azından $\omega+\omega$'ncı aşamanın,
yani $V_{\omega+\omega}$'nın bulunduğunu da kuşkusuz kabul etmeliyiz; çünkü
hiyerarşi bu kadar ilerlemeyi elbette \emph{sürdürebilir}.

Bu düşünce şöyle genelleşir: yinelemeli hiyerarşiyi kurma süreci, herhangi bir
iyi sıralama için en az o iyi sıralama kadar ilerlemelidir. Bunu yalnızca
\olref[ordinals][ordtype]{thmOrdinalRepresentation} sonucunu bir
\emph{aksiyom} sayarak güvenceye alabiliriz. Bu bize her iyi sıralamanın bir
von Neumann sıra sayısına izomorf olduğunu söyler. Her von Neumann sıra sayısı
kendi rankına eşit olacağından
\olref[ordinals][ordtype]{thmOrdinalRepresentation}, küme teorimizde ne zaman
bir iyi sıralama betimleyebilsek yinelemeli küme kurma sürecinin o iyi
sıralamanın ötesine geçmek zorunda olduğunu söyler.

Bu fikir kesinlikle \stagesinex{} ilkesinin bir sonucu gibi görünür. Ne yazık
ki amacımız bu fikirden Yerine Koyma'yı çıkarmaksa basit, teknik bir engelle
karşılaşırız: Yerine Koyma,
\olref[ordinals][ordtype]{thmOrdinalRepresentation} sonucundan kesin olarak
daha güçlüdür. (Bu gözlem \citet[\S13.2]{Potter2004} tarafından yapılmıştır;
onu \olref[sth][replacement][finiteaxiomatizability]{sec} içinde
ispatlayacağız.)

Sonuç olarak \stagesinex{} ilkesini Yerine Koyma'yı verecek biçimde
anlayacaksak ilke, hiyerarşinin herhangi bir iyi sıralamanın ötesine geçtiğini
\emph{yalnızca} söyleyemez. Bundan daha güçlü bir iddiada bulunmalıdır. Bu
amaçla \cite{Shoenfield:AST}, fikrin son derece doğal bir güçlendirmesini
önerdi: hiyerarşi hiçbir kümeyle \emph{kofinal} değildir.\footnote{G\"odel
benzer bir düşünce önermiş görünür; bkz. \citet[p.~223]{Potter2004}. G\"odel
ve \citeauthor{Shoenfield:AST} üzerine bir tartışma için bkz.
\citet[90--5]{Incurvati2020}.} Biraz daha ayrıntılı söylersek: $\tau$, kümeleri
hiyerarşinin aşamalarına gönderen bir eşlemeyse herhangi bir $A$ kümesinin
$\tau$ altındaki görüntüsü hiyerarşiyi tüketmez. Başka deyişle (şematik
olarak):
\begin{enumerate}
	\item[] \stagescofin. $A$ bir kümeyse ve her $x \in A$ için $\tau(x)$
	bir aşamaysa, her $x \in A$ için $\tau(x)$'ten sonra gelen bir aşama vardır.
\end{enumerate}
$\ZF$'nin \stagescofin{} ilkesinin uygun biçimde formelleştirilmiş bir
sürümünü ispatladığı açıktır. Tersine, \stagescofin{} ilkesinin Yerine
Koyma'yı gerekçelendirdiğini gayriresmî olarak savunabiliriz.\footnote{Bir
\stagescofin{} formelleştirmesinden yararlanarak Yerine Koyma'yı ispatlamak
daha güç olurdu; çünkü $\Z$ tek başına aşamaları tanımlayacak kadar güçlü
değildir ve dolayısıyla \stagescofin{} ilkesinin nasıl formelleştirileceği
açık değildir. Yine de bir seçenek,
\olref[sth][replacement][extrinsic]{sec} içinde tartışıldığı gibi $\LT$'nin
bir genişlemesinde çalışmaktır.} $(\forall x \in A)\lexists![y][\phi(x,y)]$
olduğunu varsayın. Her $x \in A$ için $\sigma(x)$, $\phi(x,y)$ koşulunu
sağlayan $y$; $\tau(x)$ da $\sigma(x)$'in ilk oluşturulduğu aşama olsun.
\stagescofin{} gereği $(\forall x \in A)\tau(x)\in V$ koşulunu sağlayan bir
$V$ aşaması vardır. Her $\tau(x) \in V$ ve $\sigma(x) \subseteq \tau(x)$
olduğundan Ayırma yoluyla
$\Setabs{y \in V}{(\exists x \in A)\sigma(x) = y} =
\Setabs{y}{(\exists x \in A)\phi(x,y)}$ kümesini elde edebiliriz.

\begin{prob}
	\stagescofin{} ilkesini $\ZF$ içinde formelleştirin.
\end{prob}

Dolayısıyla \stagescofin{} Yerine Koyma'yı doğrular. Ayrıca
\stagesinex{} ilkesinin \stagescofin{} ilkesini doğrulaması en azından akla
yatkındır. \stagescofin{} ilkesinin başarısız olduğunu varsayın. O zaman
hiyerarşi bazı $A$ kümeleriyle kofinaldir; yani herhangi bir $S$ aşaması için
$S \in \tau(x)$ olacak biçimde bazı $x \in A$ bulunan bir $\tau$ eşlememiz
vardır. Bu durumda hiyerarşinin sözde ``mutlak sonsuzluğunu'' kavramanın bir
yolunu elde etmiş oluruz: hiyerarşi, $\tau$'nun $A$'ya uygulanmasıyla oluşan
görüntü kümesi tarafından \emph{tüketilir}. Bu da hiyerarşinin ``mutlak
olarak sonsuz'' olduğu düşüncesini zedeler. Karşıtını alırsak:
\stagesinex{} ilkesi \stagescofin{} ilkesini gerektirir; o da Yerine
Koyma'yı gerekçelendirir.

Bu, Yerine Koyma için içsel bir gerekçelendirme sağlamaya yönelik gerçekten
umut verici bir girişimdir. Ancak nihayetinde işe yarayıp yaramadığına sizin
karar vermeniz gerekecek.

